
The effects of modifying the $W$-boson production and decay model are studied using 
event samples for the \Wlnu~process in $pp$ collisions at $\sqrt{s}=7$~TeV and
$\sqrt{s}=13$~TeV, and for $p\bar{p}$ collisions at $\sqrt{s}=1.96$~TeV.
The event generators include those used by the original experiments, along with more
recent versions with improved calculations.  The PDF sets considered include those
from the original measurements (CTEQ6M, CTEQ6.1, CTEQ6.6, CT10, NNPDF3.1, CT18,
and MSHT2020), as well as the following sets at next-to-next-to-leading order (NNLO) in
$\alpha_{\scriptsize\textrm{S}}$: NNPDF4.0~\cite{NNPDF:2021njg},
ABMP16~\cite{Alekhin:2017kpj}, CT14~\cite{Dulat:2015mca}, and MMHT2014 \cite{Harland-Lang:2014zoa}.

For the ATLAS, CDF, and D0 experiments the $m_W$ shift associated with a particular NNLO PDF set is
evaluated using the NNLO QCD calculation {\textsc{Wj-MiNNLO}} \cite{Monni:2019whf,Monni:2020nks}
implemented in {\textsc{Powheg-Box-V2}}~\cite{Nason:2004rx,Frixione:2007vw,Alioli:2010xd}.  The analysis 
is performed at the Les Houches event level~\cite{Alwall2007} without interfacing to a parton shower,
allowing for an efficient and fast processing.  The addition of the parton shower has been confirmed
to negligibly affect the shift associated with the PDF set.  A large sample is produced for each PDF set,
with weights calculated internally by \textsc{Wj-MiNNLO} in order to evaluate PDF uncertainties. 

The uncertainty associated with a given PDF set is evaluated using the next-to-leading order (NLO) QCD
calculation {\textsc{W\_ew-BMNNP}}~\cite{Barze:2012tt} implemented in the {\textsc{Powheg-Box-V2}}.  
This calculation is used for efficiency reasons and the difference in the estimated uncertainty
with respect to an NNLO calculation is expected to be negligible.
For LHCb the $m_W$ shifts and uncertainties are evaluated using
the same {\textsc{Powheg}} NLO calculation of $W$-boson production~\cite{Alioli:2008gx} as used in the original
measurement.

The modelling of the $W$-boson polarization and resonance lineshape are studied using
large \textsc{ResBos} samples corresponding to those from the Tevatron measurements:
{\textsc{ResBos-C}}~\cite{Balazs:1997xd}, used by CDF in their direct fit to the
data~\cite{CDF:2022hxs}, with an accuracy of NLO and approximate NNLL in
QCD~\cite{Isaacson:2022rts}; and {\textsc{ResBos-CP}}~\cite{Landry:2002ix}, used in the
D0 $m_W$ measurement~\cite{D0:2013jba}, with an accuracy of NNLO+NNLL in QCD.  A third sample is 
generated using {\textsc{ResBos2}}~\cite{Isaacson:2022rts} with an accuracy of NLO+NNLL in QCD and
including a full resummation of the coefficients describing the leptonic angular distributions (see
Section~\ref{sec:polarization}).  The difference between NLO and NNLO predictions of these
coefficients is studied using the DYNNLO generator~\cite{Catani:2007vq,Catani:2009sm}, which has
been confirmed to be consistent with other fixed-order calculations~\cite{comparisons}.

Electroweak corrections, primarily photon radiation in $W$-boson decay, have a large impact on the 
final-state distributions but are calculated accurately.  The experiments factorize these corrections
from PDF and QCD effects and we therefore do not include them in the sample generation.
